\ofsection{Characters}
%
\ofquote{"I am THE Basch fon Ronsenburg!"\\}{Vaan}\\
%
\includegraphics[width=\columnwidth]{./art/images/ff10-2.jpg}
%
\vfill
%
Every player creates a \accf{Character} who is a protagonist in the game world created by the GM.
To create a Level~1 character, copy or print the \accf{Character Sheet} on the next page.
It allows you track various aspects about your character and there is an example of a filled out sheet, that you can use as a guideline.
Choose your character's \accf{Name} and give a short description of him or her.
%Try to communicate with your group to properly integrate your character into the party and the game world.
%For example, your character may be part  different races or tribes exist in the world, your character may be part of one of them.
Briefly summarize your character's \accf{Story} and explain his or her motivation for joining the party, considering that this is most likely his or her first serious adventure.
Then choose a \accf{Job} for your character as explained in the following. 
Finally, the \accf{Equipment} subsection explains how you can customize your character's starting equipment.
The table below summarizes the benefits gained at further Levels, all of which are explained in detail within this section.
%
\vfill
%
\ofboxwithtitle{Example: Character Creation}
{
	We create a character named Vaan, who is a 17-year old, blonde human boy with athletic appearance. 
	Vaan is an orphan, who gets by in the big city by stealing and often acts as a father figure to other orphans.
	He dreams of owning an airship and being a sky pirate one day. 
	We communicate with the GM to make sure that Vaan fits well enough with the given setting and the rest of the party.
	From the attributes table of the job we determine Vaan's maximum HP~(20), maximum MP~(19) and AGI (4), all other attributes start at 0.
	We also note that he learns the Steal tech.
	Finally, from our initial 1500G, we buy a Mythril Knife, Myhtril Vest, a Phoenix Down and 2 Potions, which leaves us with 300G extra.
}
%
\newpage
%
\oftable{p{0.3\columnwidth} p{0.7\columnwidth}}
{\accf{Level} & \accf{Benefit Gained}}
{
	1 & Job, Beginner Equip \ofrow
	2 & Talent \ofrow
	3 & Archetype \ofrow
	4 & Limit Break, Advanced Equip \ofrow
	5 & Esper \ofrow
	6 & Custom \ofrow
	7 & Custom \ofrow
	8 & Custom, Expert Equip \ofrow
	9 & Custom \ofrow
	10 & Custom
}
%
\vfill
%
At Levels marked as \accf{Custom}, you can chose from one of the following benefits that you character gains.
You cannot pick the same benefit more than once.\ofrow
%
\ofbullet{At the start of each session, add an additional 6 to the pool of Fortune Dice.}
\ofbullet{You gain a second Esper of your choice.}
\ofbullet{You gain a second Talent of your choice.}
\ofbullet{You gain a second Limit Mode, allowing you to gain Limit Points from 2 sources. Also, the maximum Rating of your Limit Break is increased by 2.}
\ofbullet{You gain access to the second Archetype your job. You can switch between the two whenever you go to sleep and your attributes and abilities change depending on your currently active Archetype.}
\ofbullet{When you use a Spell, Tech or Reaction ability from your Archetype, its MP cost is halved and you gain 1 Limit Point.}
\ofbullet{All weapons and armor gain an additional Materia Slot as long as your have them equipped.} 
\ofbullet{You gain the ability to equip one additional weapon or armor type of your choice.} 
\ofbullet{Choose 4 from the following list of attribute bonuses. You cannot pick the same one more than twice:\newline HP+5, MP+5, STR+1, DEF+1, MAG+1, RES+1.}
%
\vfill
%
Your character's \accf{Job} determines his or her combat proficiency and progression including abilities, attributes and equipment expertise.
All available Jobs are detailed in their Job descriptions right after the character sheets.
We recommend to print or copy the description of your chosen Job to use as the second page of your character sheet.
Your character's attributes are initially 0 and increase by starting and progressing in a job.
The job description shows the attributes and abilities gained at each Level as well as the types of equipment your character can use.
When your character reaches Level~3 in his or her Job, you have to decide between one of its two \accf{Archetypes} listed on the right hand side. 
Archetypes represent different specializations of a Job and provide additional abilities and attributes.
%
\clearpage